\chapter{Wie wird das Projekt gestartet?}
\label{chap: anleitung}

Der einfachste Weg, das Projekt vollständig zum Laufen zu bringen, führt über
\emph{Docker} und \emph{Docker Compose}. Alle benötigten Dienste -- Frontend,
RAG-API, Vektordatenbank und Chat-Persistenz -- sind in einer einzigen
\texttt{docker-compose.yml} beschrieben und lassen sich mit wenigen Befehlen starten.

\section{Voraussetzungen}

Folgende Software muss auf dem Zielsystem installiert sein:

\begin{itemize}
    \item \textbf{Docker} (inkl. Docker Compose v2)
    \item Ein \textbf{OpenAI-kompatibler API-Zugang} -- entweder ein echter OpenAI API Key
    oder ein lokal laufendes Modell, das die OpenAI-API emuliert
    (z.\,B. \emph{Ollama} oder \emph{LM Studio})
\end{itemize}

\section{Konfiguration}

Vor dem ersten Start muss die Umgebungskonfiguration eingerichtet werden.
Dazu wird die mitgelieferte Beispieldatei kopiert:

\begin{lstlisting}[language=bash, caption={Konfigurationsdatei anlegen}]
cp .env.example .env
\end{lstlisting}

In der \texttt{.env}-Datei wird festgelegt, welcher LLM-Dienst verwendet werden soll.
Je nach Szenario sind unterschiedliche Einträge relevant:

\begin{itemize}
    \item \textbf{OpenAI API}: \texttt{OPENAI\_API\_KEY} auf den eigenen Key setzen,
    \texttt{OPENAI\_BASE\_URL} auf \texttt{https://api.openai.com/v1}.
    \item \textbf{Ollama (lokal)}: \texttt{OPENAI\_API\_KEY=dummy},
    \texttt{OPENAI\_BASE\_URL=http://url/}.
    \item \textbf{LM Studio}: \texttt{OPENAI\_BASE\_URL=http://host.docker.internal:1234/v1}.
\end{itemize}

\section{Erster Start: Daten laden und indizieren}

Beim allerersten Start müssen die Dokumente der FH Wedel heruntergeladen
und in die Vektordatenbank eingespeist werden. Dazu stehen zwei auskommentierte
Setup-Dienste in der \texttt{docker-compose.yml} bereit, die bei Bedarf aktiviert
werden können. Alternativ lassen sich die Skripte auch direkt ausführen:

\begin{lstlisting}[language=bash, caption={Dokumente herunterladen und indizieren}]
# PDFs von der FH-Wedel-Website herunterladen
python rag-code/download_pdfs.py

# Dokumente verarbeiten und in Qdrant einbetten
python rag-code/indexing_pipeline.py
\end{lstlisting}

Der Indizierungsschritt kann je nach Anzahl der Dokumente und verwendetem Modell
einige Minuten bis zu einer halben Stunde in Anspruch nehmen, da für jedes
Dokument Embeddings berechnet und in \emph{Qdrant} gespeichert werden.
Dieser Schritt muss nur einmalig durchgeführt werden -- die Daten bleiben
in einem Docker-Volume erhalten.

\section{Starten der Anwendung}

Sind die Dokumente indiziert, genügt ein einzelner Befehl, um alle Dienste
hochzufahren:

\begin{lstlisting}[language=bash, caption={Alle Dienste starten}]
docker compose up frontend -d
\end{lstlisting}

Durch die Abhängigkeitsstruktur in der \texttt{docker-compose.yml} werden dabei
automatisch auch alle anderen benötigten Dienste gestartet. Nach kurzer Initialisierungszeit
sind die Dienste über folgende Adressen erreichbar:

\begin{itemize}
    \item \textbf{Frontend} (Next.js-Weboberfläche): \texttt{http://localhost:3000}
    \item \textbf{RAG-API} (FastAPI-Backend): \texttt{http://localhost:8000}
    \item \textbf{Qdrant-Dashboard}: \texttt{http://localhost:6333/dashboard}
    \item \textbf{MongoDB}: \texttt{localhost:27017}
\end{itemize}

Die Weboberfläche unter Port 3000 ist der primäre Einstiegspunkt für den Endnutzer.
Dort kann nach der Auswahl eines Studiengangs direkt mit dem Chatbot interagiert werden.

\section{Nutzung der CLI-Schnittstelle}

Alternativ zur Weboberfläche steht eine interaktive Kommandozeilenschnittstelle
zur Verfügung. Diese ist besonders während der Entwicklung oder für schnelle
Tests praktisch:

\begin{lstlisting}[language=bash, caption={CLI-Modus starten}]
docker compose run --rm rag-cli
\end{lstlisting}

\section{Logs und Fehlerdiagnose}

Die Logs aller laufenden Dienste lassen sich mit folgendem Befehl verfolgen:

\begin{lstlisting}[language=bash, caption={Logs beobachten}]
docker compose logs -f
\end{lstlisting}

Zum Stoppen aller Dienste:

\begin{lstlisting}[language=bash, caption={Alle Dienste stoppen}]
docker compose down
\end{lstlisting}

\section{Hinweise zur Entwicklung}

Für die aktive Entwicklung am Projekt sind in der \texttt{docker-compose.yml}
Bind-Mounts eingerichtet, sodass Änderungen am Quellcode im laufenden Container
direkt sichtbar sind. Das Backend startet mit dem Flag \texttt{--reload}, was
bei Dateiänderungen einen automatischen Neustart des \emph{uvicorn}-Servers
auslöst. Gleiches gilt für das Next.js-Frontend, das im Entwicklungsmodus
ebenfalls mit Hot Reload betrieben wird.
